\documentclass[12pt, a4paper]{article}

% ── PACKAGES ──────────────────────────────────────────────────
\usepackage[utf8]{inputenc}
\usepackage[T1]{fontenc}
\usepackage{geometry}
\geometry{margin=2.5cm}
\usepackage{hyperref}
\hypersetup{
    colorlinks=true,
    linkcolor=blue,
    urlcolor=blue,
    citecolor=blue,
    pdftitle={Attractor Geometry of Life: Pre-Registration of
              Geometric Derivations, Experimental Predictions,
              and Falsifiability Conditions},
    pdfauthor={Eric Robert Lawson},
    pdfkeywords={attractor geometry, LECA, developmental arrest,
                 Reproduction Theorem, de-extinction, LUCA,
                 ribosome, triadic invariant, Dickinsonia,
                 cross-kingdom convergence, evolutionary biology,
                 OrganismCore}
}
\usepackage{amsmath}
\usepackage{amssymb}
\usepackage{amsthm}
\usepackage{booktabs}
\usepackage{longtable}
\usepackage{array}
\usepackage{parskip}
\usepackage{microtype}
\usepackage{authblk}
\usepackage{titlesec}
\usepackage{fancyhdr}
\usepackage{xcolor}
\usepackage{float}
\usepackage{graphicx}
\usepackage{setspace}
\usepackage{enumitem}
\usepackage{tabularx}
\usepackage{multirow}
\usepackage{doi}

% ── COLOURS ───────────────────────────────────────────────────
\definecolor{repobox}{RGB}{240,247,255}
\definecolor{lockbox}{RGB}{245,255,245}
\definecolor{warnbox}{RGB}{255,248,220}
\definecolor{theorembox}{RGB}{248,240,255}

% ── THEOREM ENVIRONMENTS ──────────────────────────────────────
\newtheorem{theorem}{Theorem}[section]
\newtheorem{corollary}[theorem]{Corollary}
\newtheorem{lemma}[theorem]{Lemma}
\theoremstyle{definition}
\newtheorem{definition}{Definition}[section]
\newtheorem{prediction}{Prediction}[section]
\theoremstyle{remark}
\newtheorem{remark}{Remark}[section]

% ── HEADER / FOOTER ───────────────────────────────────────────
\pagestyle{fancy}
\fancyhf{}
\rhead{\small OrganismCore \textbar{} 2026-03-12}
\lhead{\small Attractor Geometry of Life --- Pre-Registration}
\rfoot{\small Page \thepage}
\renewcommand{\headrulewidth}{0.4pt}

% ── SECTION FORMATTING ────────────────────────────────────────
\titleformat{\section}
  {\large\bfseries}{\thesection}{1em}{}[\titlerule]
\titleformat{\subsection}
  {\normalsize\bfseries}{\thesubsection}{1em}{}

% ══════════════════════════════════════════════════════════════
% TITLE
% ══════════════════════════════════════════════════════════════
\title{
    \vspace{-1em}
    \textbf{Attractor Geometry of Life:}\\[0.5em]
    \large\textbf{Pre-Registration of Geometric Derivations,
    Experimental Predictions,}\\
    \large\textbf{and Falsifiability Conditions}\\[0.6em]
    \normalsize\textit{Covering the Reproduction Theorem,
    the Inverse Accessibility Theorem,}\\
    \normalsize\textit{the LECA Resurrection Protocol,
    the Dickinsonia Protocol,}\\
    \normalsize\textit{the Cross-Kingdom Convergence Prediction,}\\
    \normalsize\textit{and the Geometric Status of the LUCA}\\[0.6em]
    \small Derived from Waddington Landscape Attractor
    Geometry and the Triadic Invariant\\
    \small OrganismCore Document:
    ATTRACTOR-LIFE-PREREGISTRATION-v1.0
}

\author[1]{Eric Robert Lawson}
\affil[1]{
    OrganismCore\\
    \href{mailto:OrganismCore@proton.me}
         {OrganismCore@proton.me}\\
    ORCID:
    \href{https://orcid.org/0009-0002-0414-6544}
         {0009-0002-0414-6544}\\
    Repository:
    \href{https://github.com/Eric-Robert-Lawson/attractor-oncology}
         {github.com/Eric-Robert-Lawson/attractor-oncology}
}

\date{
    Pre-registration date: 2026-03-12\\
    \small\textit{All predictions in this document are
    pre-specified prior to any experimental confirmation.}\\
    \small\textit{This document constitutes the formal
    pre-registration record.}
}

% ══════════════════════════════════════════════════════════════
\begin{document}
% ══════════════════════════════════════════════════════════════

\maketitle
\thispagestyle{fancy}

% ── REPOSITORY POINTER ────────────────────────────────────────
\begin{center}
\colorbox{repobox}{
\begin{minipage}{0.88\textwidth}
\vspace{0.5em}
\centering
\textbf{Full derivation record, protocols, reasoning
artifacts, and pre-registration documents:}\\[0.4em]
\href{https://github.com/Eric-Robert-Lawson/attractor-oncology}
{\texttt{https://github.com/Eric-Robert-Lawson/attractor-oncology}}
\\[0.3em]
\small
Repository ID: 1172048282 \textbar{}
License: MIT \textbar{}
Language: Python \textbar{}
Visibility: Public\\[0.2em]
Topics: attractor \textperiodcentered{}
        bioinformatics \textperiodcentered{}
        cancer-genomics \textperiodcentered{}
        computational-oncology \textperiodcentered{}
        drug-discovery \textperiodcentered{}
        epigenetics \textperiodcentered{}
        evolutionary-biology \textperiodcentered{}
        developmental-biology \textperiodcentered{}
        de-extinction\\[0.3em]
All derivations are timestamped and version-controlled.
All predictions are pre-specified before experimental
confirmation.\\
\vspace{0.5em}
\end{minipage}
}
\end{center}

\vspace{0.8em}

% ── ABSTRACT ──────────────────────────────────────────────────
\begin{abstract}
\noindent
We present the formal pre-registration record for a set of
geometric derivations, theorems, and experimentally testable
predictions produced by the Attractor Geometry framework ---
a principles-first geometric framework derived from the
triadic invariant as applied to biological identity at all
scales. The framework was originally developed to derive
cancer drug targets from epigenetic geometry and has been
extended to evolutionary biology, developmental biology, and
synthetic organism construction.

The central result is the \textbf{Reproduction Theorem}: the
developmental trajectory of any organism from informational
encoding to adult form is geometrically identical to the
evolutionary trajectory of its lineage, compressed into
developmental time. This is established by proof of
elimination: no alternative coherent trajectory from
informational biological encoding to a living organism has
been described in 150 years of developmental biology.

The theorem generates the \textbf{Inverse Accessibility
Theorem}: the oldest recoverable organism is the easiest and
cheapest to recover, because evolutionary age is inversely
proportional to developmental position, and the earlier the
developmental position the simpler the arrest intervention.

The most accessible experimental consequence is the
\textbf{LECA Resurrection Protocol}: developmental arrest of
any eukaryotic cell before its first multicellular
organisational step will produce an organism morphologically
consistent with the Last Eukaryotic Common Ancestor
($\sim$1.5\,Ga). This result is predicted to be
morphologically identical whether the starting material is an
animal embryo, a plant seed, or a fungal spore --- the
\textbf{Cross-Kingdom Convergence Prediction} --- while being
genomically distinct under sequencing. The experiment costs
\$200--400 per run and is executable in any standard cell
biology laboratory in days.

We further establish the geometric status of the LUCA: the
ribosome's peptidyl transferase center constitutes the
continuously executing Row~3 identity of the Last Universal
Common Ancestor. The LUCA is not extinct by the precise
definition of the triadic invariant. The ribosome is
identified as the topological origin of the structural space
of all life and as the medium on which every developmental
arrest runs --- present inside every recovered ancestral
organism, older than every attractor it will ever be found in.

All predictions are pre-specified with explicit falsifiability
conditions. Pre-registration timestamp: 2026-03-12.
\end{abstract}

\vspace{0.5em}
\noindent\textbf{Keywords:}
attractor geometry; triadic invariant; Reproduction Theorem;
Inverse Accessibility Theorem; LECA; LUCA; ribosome;
developmental arrest; de-extinction; Dickinsonia; Ediacaran;
cross-kingdom convergence; Waddington landscape; OrganismCore

\newpage
\tableofcontents
\newpage

% ═══════════════════════��══════════════════════════════════════
\section{Document Metadata}
% ══════════════════════════════════���═══════════════════════════

\begin{center}
\begin{tabular}{ll}
\toprule
\textbf{Field} & \textbf{Value} \\
\midrule
Document ID       & ATTRACTOR-LIFE-PREREGISTRATION-v1.0 \\
Series            & Attractor Geometry of Life \\
Type              & Pre-Registration Record \\
Date              & 2026-03-12 \\
Author            & Eric Robert Lawson \\
Organisation      & OrganismCore \\
ORCID             &
  \href{https://orcid.org/0009-0002-0414-6544}
       {0009-0002-0414-6544} \\
Contact           &
  \href{mailto:OrganismCore@proton.me}
       {OrganismCore@proton.me} \\
Status            & Permanent pre-registration record \\
Repository        &
  \href{https://github.com/Eric-Robert-Lawson/attractor-oncology}
       {Eric-Robert-Lawson/attractor-oncology} \\
Repository ID     & 1172048282 \\
License           & MIT \\
Biological contradictions & 0 \\
Type B failures   & 0 \\
Pre-registration timestamp & 2026-03-12 \\
Prior framework   & FOXA1-EZH2-RATIO-MASTER-RA (2026-03-06) \\
\bottomrule
\end{tabular}
\end{center}

% ══════════════════════════════════════════════════════════════
\section{Preamble: The Derivation Path}
% ══════════════════════════════════════════════════════════════

This document records geometric derivations produced by Eric
Robert Lawson, an independent autodidactic researcher with no
institutional affiliation, no laboratory access, and no formal
training in biology, working from first principles using the
triadic invariant as the foundational structure.

The derivation path was not planned. It followed geometrically:

\begin{enumerate}[label=\arabic*.]
  \item Cancer drug target geometry $\rightarrow$ the
        Identity Axis Theorem: the ratio $R = A/H$
        (FOXA1/EZH2) as the epigenetic identity axis of
        committed breast epithelial cell types. Validated
        across $\sim$7,500 patients in seven independent
        datasets on four measurement platforms
        (see FOXA1-EZH2-RATIO-MASTER-RA, 2026-03-06).

  \item Identity axis geometry $\rightarrow$ the triadic
        invariant as the causal structure of all biological
        identity at every scale.

  \item Triadic invariant applied to organism persistence
        across time $\rightarrow$ the Reproduction Theorem.

  \item Reproduction Theorem $\rightarrow$ Inverse
        Accessibility Theorem, LECA protocol, Dickinsonia
        protocol, plant Ediacaran protocol.

  \item Ancestral access map $\rightarrow$ geometric status
        of the LUCA: the ribosome as the LUCA's continuously
        executing Row~3 identity. The LUCA is not extinct.

  \item Ribosome as LUCA $\rightarrow$ the ribosome as
        topological origin of the structural space of all
        life and as the medium on which every developmental
        arrest runs.
\end{enumerate}

The derivation of the ribosome's geometric role was produced
without prior knowledge of ribosome structure or function.
The structure was identified as the ribosome only after the
geometric role was derived. The literature confirmed the
observation post-derivation. This epistemic pattern ---
geometry first, identification second, literature confirmation
third --- is consistent across all results in this framework.

% ══════════════════════════════════════════════════════════════
\section{The Triadic Invariant: Foundational Structure}
% ══════════════════════════════════════════════════════════════

\begin{definition}[The Triadic Invariant]
All stable biological identity at every scale is constituted
by a three-row structure:
\begin{itemize}
  \item \textbf{Row 1} (Informational Basis): The substrate
    encoding the attractor geometry --- the compressed record
    of evolutionary navigation.
  \item \textbf{Row 2} (Gap Navigation): The process by which
    Row~1 is decompressed into a physical traversal of the
    attractor landscape.
  \item \textbf{Row 3} (Stable Product): The coherent,
    persistent, self-referential identity that is the stable
    endpoint of Row~2 navigation and that maintains the
    conditions of its own production.
\end{itemize}
\end{definition}

\begin{definition}[Extinction]
An organism is extinct when its Row~3 product --- its
coherent, persistent, self-maintaining identity --- ceases
to be produced.
\end{definition}

\begin{definition}[Attractor]
An attractor is a stable configuration in the epigenetic or
developmental landscape that Row~2 navigation reaches and
maintains. Attractors are sequentially dependent: each
attractor requires the stable product of the preceding
attractor as its Row~1 substrate.
\end{definition}

\subsection{Prior validation of the triadic invariant}

The triadic invariant was first validated in the cancer
framework. Applied to breast cancer, it produced the
FOXA1/EZH2 ratio as the epigenetic identity axis --- the
continuous measure of the balance between luminal identity
maintenance (FOXA1) and identity suppression (EZH2). This
prediction was locked before confirmatory analysis and
confirmed across $\sim$7,500 patients in seven independent
datasets with zero biological contradictions
(see repository: \texttt{Cancer\_Research/BRCA/DEEP\_DIVE/\\
MAJOR\_FINDINGS/FOXA1\_EZH2\_RATIO/FOXA1\_EZH2\_BCRA\_Ratio.tex}).

The same invariant, applied to the problem of organism
persistence across evolutionary time, generates all results
in this document.

% ══════════════════════════════════════════════════════════════
\section{The Reproduction Theorem}
% ══════════════════════════════════════════════════════════════

\subsection{Statement}

\begin{theorem}[The Reproduction Theorem]
The developmental trajectory of any organism from
informational encoding (fertilised egg or germinating seed)
to adult organism is geometrically identical to the
evolutionary trajectory of its lineage, expressed at a
compressed temporal scale.

They are not analogous. They are not parallel. They are the
same sequence of attractor states, in the same order, from
the same informational starting point, to the same terminal
state.
\end{theorem}

\subsection{The developmental-evolutionary correspondence}

\begin{center}
\begin{tabular}{p{4.5cm}p{4.5cm}p{2.5cm}}
\toprule
\textbf{Developmental Stage} &
\textbf{Evolutionary Grade} &
\textbf{Age (Ma)} \\
\midrule
Fertilised egg / first cleavages
  & First multicellular aggregates & $\sim$800 \\
Blastula
  & Colonial multicellular, earliest animal grade & $\sim$700 \\
Gastrula
  & Diploblast / early bilaterian & $\sim$560--600 \\
Neurula
  & Urbilaterian / early deuterostome & $\sim$540--560 \\
Pharyngeal arch stage (phylotypic)
  & Vertebrate ancestor & $\sim$500 \\
Organogenesis
  & Specific vertebrate lineage grades & $\sim$400--300 \\
Late fetal
  & Mammalian ancestor & $\sim$200 \\
Birth / adult
  & Species-specific & Present \\
\bottomrule
\end{tabular}
\end{center}

\noindent This correspondence is confirmed by the
developmental hourglass model across all major animal phyla
\cite{kalinka2010, irie2011} and extended to plants
\cite{plant_hourglass2024}.

\subsection{Proof by elimination of alternatives}

\begin{proof}
Let $\mathcal{T}_D$ denote the developmental trajectory and
$\mathcal{T}_E$ the evolutionary trajectory of lineage $L$.

\textbf{Premise 1.} There exists exactly one known trajectory
$\mathcal{T}_D$ from informational encoding to a living
organism of any species. No alternative trajectory has been
observed or demonstrated.

\textbf{Premise 2.} The stages of $\mathcal{T}_D$ are
\emph{sequentially dependent attractors}: each stage $A_i$
is only accessible from $A_{i-1}$. Disruption of any stage
is lethal. The order is geometrically necessary, not
conventional.

\textbf{Premise 3.} The stages of $\mathcal{T}_D$ correspond
stage-by-stage to the grades occupied by the lineage during
$\mathcal{T}_E$, from earliest to most recently derived
(confirmed by the developmental hourglass across all major
animal and plant phyla).

\textbf{Claim.} $\mathcal{T}_D = \mathcal{T}_E$ (compressed).

\textbf{Proof.} Suppose $\mathcal{T}_D \neq \mathcal{T}_E$.
Then one of the following must hold:

\begin{description}
  \item[Alt.\ A.] $\mathcal{T}_D$ skips one or more attractor
    states in $\mathcal{T}_E$. Excluded by Premise~2:
    skipping any attractor is lethal.
  \item[Alt.\ B.] $\mathcal{T}_D$ traverses attractor states
    in a different order from $\mathcal{T}_E$. Excluded by
    sequential dependence: each stage's molecular preconditions
    are established only by the preceding stage in the
    observed order.
  \item[Alt.\ C.] $\mathcal{T}_D$ traverses different
    attractor states and still reaches the same terminal
    organism. Requires the terminal attractor to be reachable
    from a sequence no lineage ever traversed. Not demonstrated
    in 150 years of developmental biology.
\end{description}

All alternatives are excluded.
Therefore $\mathcal{T}_D = \mathcal{T}_E$ (compressed).
$\square$
\end{proof}

\subsection{Corollaries}

\begin{corollary}
Every reproductive cycle re-executes the complete
evolutionary history of the organism's lineage in compressed
time.
\end{corollary}

\begin{corollary}[Developmental Arrest as Evolutionary Access]
Arresting the developmental program of any organism at stage
$A_i$ produces the organism that occupied attractor $A_i$ in
the evolutionary history of the lineage --- not a
morphological analog, but the attractor itself, stabilised
rather than passed through.
\end{corollary}

\begin{corollary}[Lineage Specificity Constraint]
Only attractors that the starting organism's lineage actually
traversed are accessible by developmental arrest of that
organism. Attractors from a different lineage are not encoded
in the developmental program and cannot be accessed.
\end{corollary}

\subsection{Literature convergences confirming the theorem}

The following independent empirical results confirm the
theorem's geometric structure. In each case the result was
found by the existing literature without knowledge of the
framework.

\begin{center}
\begin{tabular}{p{5.5cm}p{6cm}}
\toprule
\textbf{Framework derivation} &
\textbf{Independent empirical confirmation} \\
\midrule
Developmental hourglass geometrically forced
  & Kalinka et al.\ 2010, \textit{Nature}:
    hourglass confirmed across \textit{Drosophila} species \\
Phylotypic stage conserved across all vertebrates
  & Irie \& Kuratani 2011, \textit{Nat.\ Commun.}:
    vertebrate phylotypic stage confirmed transcriptomically \\
Hourglass is geometric necessity, not contingency
  & PLOS Comp.\ Biol.\ 2023: hourglass emerges necessarily
    from numerical evolution simulations \\
Plant developmental program compresses plant
evolutionary history
  & Nat.\ Commun.\ 2024: plant embryonic hourglass
    confirmed \cite{plant_hourglass2024} \\
Forward lock prevents developmental regression
  & \textit{Nature} 2024: LINE-1 RNA prevents
    developmental regression \cite{line1_2024} \\
\bottomrule
\end{tabular}
\end{center}

% ══════════════════════════════════════════════════════════════
\section{The Inverse Accessibility Theorem}
% ══════════════════════════════════════════════════════════════

\begin{theorem}[The Inverse Accessibility Theorem]
Under the Reproduction Theorem, the accessibility of an
ancestral attractor by developmental arrest is inversely
proportional to its evolutionary age:

\[
  \mathrm{Accessibility} \;\propto\;
  \frac{1}{\mathrm{Evolutionary\;Age}}
\]

The older the target organism: the earlier it sits in the
developmental program; the simpler the arrest intervention;
the lower the cost; the shorter the experimental timeline.
\end{theorem}

\begin{proof}
By the Reproduction Theorem, $\mathcal{T}_D =
\mathcal{T}_E$ (compressed). Stages are ordered from most
ancient ($A_1$) to most recently derived ($A_n$). Earlier
stages require simpler intervention because: (1)~the
process to be arrested has not yet begun (prevention rather
than midstream arrest); (2)~the resulting organism requires
less environmental maintenance at a lower organisational
grade. Therefore evolutionary age $\propto$ developmental
earliness $\propto$ simplicity. $\square$
\end{proof}

\begin{table}[H]
\centering
\caption{Inverse Accessibility Theorem: empirical evidence
across de-extinction approaches.}
\label{tab:inverse}
\begin{tabularx}{\textwidth}{lXXll}
\toprule
\textbf{Target} & \textbf{Age} &
\textbf{Starting material} &
\textbf{Cost/run} & \textbf{Timeline} \\
\midrule
Mammoth & 4,000 yr
  & Ancient DNA required
  & \$225M (not done) & 10$+$ yr \\
Thylacine & 86 yr
  & Preserved tissue
  & N/A (not done) & 10 yr \\
Non-avian dinosaur & 66 Ma
  & Bird embryo
  & N/A & 10--20 yr \\
Dickinsonia & $\sim$558 Ma
  & Any animal
  & \$550/run & 5--8 days \\
Plant Ediacaran & $\sim$700 Ma
  & Any plant/seed
  & \$300--500/run & Days \\
\textbf{LECA} & \textbf{$\sim$1,500 Ma}
  & \textbf{Any eukaryote}
  & \textbf{\$200--400/run} & \textbf{Days} \\
\bottomrule
\end{tabularx}
\end{table}

\noindent The oldest recoverable organism is the cheapest
and fastest to recover. This is a geometric consequence of
the Reproduction Theorem, not a contingent empirical finding.

% ══════════════════════════════════════════════════════════════
\section{The Complete Ancestral Access Map}
% ══════════════════════════════════════════════════════════════

\begin{table}[H]
\centering
\caption{Complete ancestral access map for eukaryotic life.}
\label{tab:access}
\begin{tabularx}{\textwidth}{lXlX}
\toprule
\textbf{Grade} & \textbf{Age} &
\textbf{Status} & \textbf{Access route} \\
\midrule
Pre-cellular life & $>$4.0 Ga
  & \textcolor{red}{Inaccessible}
  & No genome existed. The wall. \\
LUCA & 3.8--4.0 Ga
  & \textcolor{blue}{Already present}
  & Not extinct. Ribosome in every cell. \\
\textbf{LECA} & \textbf{$\sim$1.5 Ga}
  & \textbf{\textcolor{green!60!black}{Accessible}}
  & \textbf{Any eukaryote. The floor.} \\
Plant Ediacaran & $\sim$700 Ma
  & \textcolor{green!60!black}{Accessible}
  & Any plant/seed. \\
Animal Ediacaran (Dickinsonia) & $\sim$540--600 Ma
  & \textcolor{green!60!black}{Accessible}
  & Any animal. \\
All modern grades & 0
  & Present
  & No arrest needed. \\
\bottomrule
\end{tabularx}
\end{table}

\subsection{The wall}

The developmental arrest method cannot access pre-cellular
life. The Reproduction Theorem requires a genome as its
encoding medium. Pre-cellular chemistry predates the genome.
No developmental program re-executes it. This is a
structural consequence of the theorem, not a technical
limitation.

\subsection{The floor}

The LECA is the floor of the developmental arrest method.
It sits below every eukaryotic lineage divergence --- the
plant-animal divergence ($\sim$1.5 Ga), the animal-fungal
divergence ($\sim$1.0 Ga) --- and is therefore accessible
from every eukaryotic lineage. No other ancestral attractor
is accessible from all three kingdoms simultaneously.

% ══════════════════════════════════════════════════════════════
\section{The LECA Resurrection Protocol}
% ══════════════════════════════════════════════════════════════

\subsection{The target organism: derived characteristics}

Based on conserved gene presence across all eukaryotic
kingdoms \cite{luca2024}, the LECA's characteristics are
derivable without ancient genetic material:

\begin{itemize}[noitemsep]
  \item Single cell, $\sim$10--20\,$\mu$m diameter
  \item Nucleus with nuclear envelope and nuclear pore
        complexes
  \item Mitochondria: already integrated via endosymbiosis
  \item Cytoskeletal network: actin and tubulin both present
  \item One or more flagella for motility
  \item Heterotrophic: requires external organic nutrient
  \item Aquatic: freshwater or shallow marine environment
  \item Full meiotic machinery conserved:
        Spo11, DMC1, RAD51, cohesins
  \item Genome: estimated 4,000--6,000 protein-coding genes
  \item No multicellular organisation of any kind
\end{itemize}

\subsection{The arrest principle}

To access the LECA attractor, the developmental program must
be arrested before the first multicellular organisational
step:

\begin{itemize}
  \item \textbf{Animal (mammalian):} Before compaction
    (morula stage). Target: 1--4 cell stage. The cell is in
    the LECA-grade single-cell attractor. Arrest using CDK1
    inhibitor (RO-3306, $\sim$\$80) or aphidicolin.

  \item \textbf{Plant:} Zygote before first asymmetric
    division. Inhibit the asymmetric division programme
    using auxin transport inhibitor NPA. The first division
    IS the first step away from the LECA attractor.

  \item \textbf{Fungal spore (recommended first attempt):}
    The spore before germination pore formation. Withhold
    nitrogen source; maintain below germination temperature.
    No animal facility required. No ethics committee approval
    required. Any molecular biology laboratory with yeast
    experience can execute this protocol immediately.
\end{itemize}

\subsection{Protocol: fungal version}

\textbf{Starting material:}
\textit{Saccharomyces cerevisiae} spores or
\textit{Aspergillus niger} conidia. Commercially available.

\textbf{Arrest intervention:}
\begin{enumerate}[label=\arabic*.]
  \item Suspend spores in minimal aqueous medium.
  \item Withhold nitrogen source (omit amino acids and
        ammonium sulphate).
  \item Temperature: 20--22$^\circ$C (below optimal
        germination threshold).
  \item Confirm arrest at 2h, 6h, 24h by light microscopy:
        absence of germ tube or hyphal emergence.
\end{enumerate}

\textbf{Medium composition:}
\begin{center}
\begin{tabular}{ll}
\toprule
\textbf{Component} & \textbf{Concentration} \\
\midrule
Distilled water & Base \\
NaCl & 10 mM \\
KCl & 5 mM \\
MgCl$_2$ & 2 mM \\
CaCl$_2$ & 1 mM \\
FeSO$_4$ & 1 mM \\
KH$_2$PO$_4$ & 1 mM \\
Glucose & 0.1\% (heterotrophic maintenance only) \\
pH & 6.5--7.0 \\
Atmosphere & Microaerophilic preferred \\
\bottomrule
\end{tabular}
\end{center}

\textbf{Validation criterion (morphological):}
Single-cell form maintained; nucleus visible
($\geq$40$\times$ objective); diameter 10--30\,$\mu$m;
no multicellular adhesion events; single-cell state
maintained for $\geq$72 hours without progression toward
germination.

\textbf{Estimated cost:} \$200--400 total.
Equipment: biosafety cabinet, incubator, light microscope.
All standard in any cell biology laboratory.

\textbf{Timeline to first result:} 24--72 hours.

% ══════════════════════════════════════════════════════════════
\section{The Dickinsonia Protocol}
% ══════════════════════════════════════════════════════════════

\subsection{The target organism}

\textit{Dickinsonia} was an Ediacaran organism ($\sim$558 Ma)
confirmed as an early animal by cholesterol biomarker
analysis \cite{bobrovskiy2018}. It represents the Ediacaran
organisational grade: multicellular, bilateral or
near-bilateral symmetry, 2--6 cell types, no organ systems,
no dedicated nervous or circulatory system, modular growth
by accretion.

The Ediacaran grade is not lost. It is running in every
animal developmental program, in every reproductive cycle,
at the blastula-to-early-gastrula transition. It has been
running continuously for 540 million years.

\subsection{The arrest principle}

\textbf{Recommended starting material:}
\textit{Trichoplax adhaerens} --- the only living animal at
the base of the animal tree, already occupying an attractor
adjacent to the Ediacaran grade.

\textbf{Intervention:}
\begin{enumerate}[label=\arabic*.]
  \item Activate Wnt/$\beta$-catenin signalling at the
        blastula-to-gastrula transition to specify the
        primary axis.
  \item Inhibit further mesoderm specification programmes
        (BMP gradient manipulation).
  \item Maintain in low-oxygen ($<$1\% O$_2$) aquatic
        conditions approximating Ediacaran shallow marine
        environment.
  \item Allow the Ediacaran-grade multicellular attractor
        to stabilise.
\end{enumerate}

\textbf{Estimated cost:} \$550 per run.
\textbf{Timeline:} 5--8 days to stable multicellular form.

% ══════════════════════════════════════════════════════════════
\section{Pre-Registered Predictions}
% ══════════════════════════════════════════════════════════════

\noindent All predictions below are pre-specified as of
2026-03-12. Each carries its geometric basis, measurable
output, and explicit failure condition.

\subsection{LECA predictions}

\begin{prediction}[Morphological Consistency]
Developmental arrest of any eukaryotic organism before
multicellular organisational commitment will produce a cell
morphologically consistent with the LECA-grade
characteristics specified in Section~7.1.

\textbf{Measurable output:} Single cell; 10--30\,$\mu$m;
nucleus visible; motile-capable; maintained in single-cell
state for $\geq$72h without differentiation.

\textbf{Failure condition:} The arrested cell is not viable,
or immediately progresses to multicellular organisation
despite the arrest intervention, or shows morphology
inconsistent with LECA-grade characteristics.
\end{prediction}

\begin{prediction}[Cross-Kingdom Morphological Convergence]
\label{pred:ckc}
Developmental arrest at the LECA grade from:
(1)~a mammalian embryo (e.g.\ mouse zygote),
(2)~a plant seed (e.g.\ \textit{Arabidopsis thaliana}), and
(3)~a fungal spore (e.g.\ \textit{S.\ cerevisiae})
will produce organisms that are \textbf{morphologically
indistinguishable} under light microscopy.

\textbf{Measurable output:} A blinded panel of developmental
biologists given unlabelled samples cannot identify the
starting kingdom above chance ($> 1/3$ correct) from
morphological examination alone.

\textbf{Geometric basis:} The active genes at the LECA grade
are $\sim$90--95\% identical across all eukaryotic kingdoms.
The lineage-specific genes that differentiate animal, plant,
and fungal development are not yet active at this arrest
point. The morphology is determined by the active genes.
Therefore the morphology converges.

\textbf{Failure condition:} The arrested organisms show
systematically different morphologies attributable to
lineage-specific programs such that a blinded observer can
identify the starting kingdom above chance.
\end{prediction}

\begin{prediction}[Genomic Distinctness Under Morphological
Identity]
The morphologically identical arrested organisms of
Prediction~\ref{pred:ckc} will be genomically distinct:
whole-genome sequencing will correctly identify the starting
kingdom with high confidence.

\textbf{Significance:} Morphologically identical $+$
genomically distinct $=$ the attractor is real and
independent of the genomic substrate. The geometry produces
the same form from different starting material.

\textbf{Failure condition:} Sequencing cannot distinguish
starting kingdoms above chance.
\end{prediction}

\begin{prediction}[Residual Genome Silencing Profile]
Transcriptomic analysis (RNA-seq) of arrested LECA-grade
organisms will show:
\begin{enumerate}[label=(\alph*)]
  \item High expression of LECA-grade conserved ancestral
        genes (Wnt components, core cell cycle machinery,
        cytoskeletal genes, mitochondrial genes).
  \item Near-zero expression of lineage-specific
        post-LECA developmental genes.
  \item Epigenetic silencing marks (H3K27me3, DNA
        methylation) on silent lineage-specific regions.
\end{enumerate}

\textbf{Failure condition:} Lineage-specific post-LECA
genes are actively expressed at levels comparable to
normally developing embryos at the same stage.
\end{prediction}

\begin{prediction}[CRISPR Excision Stabilises the Attractor]
Excision of post-LECA developmental programmes from an
arrested organism using CRISPR will increase attractor
stability relative to non-excised controls:
\begin{enumerate}[label=(\alph*)]
  \item Longer duration of LECA-grade morphological
        maintenance.
  \item Reduced frequency of spontaneous differentiation.
  \item Reduced metabolic cost (elimination of epigenetic
        maintenance burden for silenced regions).
\end{enumerate}

\textbf{Geometric basis:} Excision removes the silenced
downstream genome, eliminating epigenetic landscape tension
from heterochromatin adjacent to active LECA-grade genes.
The LECA attractor becomes the only attractor in the
organism's epigenetic landscape. Analogous to Hub removal
in the cancer framework collapsing the cancer attractor.

\textbf{Failure condition:} No measurable difference in
stability between excised and non-excised arrested organisms.
\end{prediction}

\subsection{Dickinsonia / Ediacaran predictions}

\begin{prediction}[Ediacaran-Grade Organism from Animal Arrest]
Developmental arrest of any animal embryo at the
blastula-to-early-gastrula transition, with Wnt axis
activation and inhibition of further mesoderm specification,
in low-oxygen aquatic conditions, will produce a
multicellular organism morphologically consistent with the
Ediacaran organisational grade: bilateral or near-bilateral
symmetry; 2--6 cell types; no organ systems; modular
accretion growth.

\textbf{Recommended starting material:}
\textit{Trichoplax adhaerens}.

\textbf{Estimated cost:} \$550 per run.

\textbf{Failure condition:} The arrested organism does not
develop a stable multicellular form, or develops organ
system complexity inconsistent with the Ediacaran grade,
with no structural explanation.
\end{prediction}

\begin{prediction}[Cross-Species Animal Convergence at
Ediacaran Grade]
Developmental arrest at the Ediacaran grade from
phylogenetically distant starting animal species (minimum:
one vertebrate, one arthropod, one cnidarian or placozoan)
will produce organisms morphologically indistinguishable
under light microscopy.

\textbf{Failure condition:} Arrested organisms from
different animal lineages show systematically different
morphologies attributable to lineage-specific programs
expressed at the arrest point.
\end{prediction}

\subsection{LINE-1 independent access prediction}

\begin{prediction}[LINE-1 Inhibition as Independent Access
Pathway]
Inhibition of LINE-1 retrotransposon activity in any animal
embryo at or before the Ediacaran-grade developmental stage
will produce spontaneous developmental regression toward the
Ediacaran attractor without deliberate signalling
intervention.

\textbf{Mechanistic basis:} The 2024 discovery that LINE-1
RNA functions as a forward-lock preventing developmental
regression \cite{line1_2024}. Inhibiting this lock should
produce spontaneous regression.

\textbf{Significance:} If both the Wnt-activation pathway
and the LINE-1-inhibition pathway produce organisms at the
same attractor from the same or different starting species,
this constitutes convergence from two mechanistically
independent directions onto the same attractor. This is the
most powerful possible geometric confirmation.

\textbf{Failure condition:} LINE-1 inhibition does not
produce developmental regression, or produces regression to
an attractor other than the Ediacaran grade.
\end{prediction}

% ══════════════════════════════════════════════════════════════
\section{The Geometric Status of the LUCA}
% ══════════════════════════════════════════════════════════════

\subsection{The LUCA is not extinct}

\begin{theorem}[The LUCA Is Not Extinct]
The Last Universal Common Ancestor is not extinct.

By Definition~2, extinction requires cessation of the
organism's Row~3 product. The LUCA's Row~3 product ---
template-driven peptide bond formation by the peptidyl
transferase centre (PTC) of the ribosome --- has not
ceased. It has been executed continuously in every living
cell on Earth since the LUCA existed. Therefore the LUCA
is not extinct.
\end{theorem}

\begin{proof}
The PTC of the large ribosomal subunit is structurally and
sequentially identical across all domains of life: Bacteria,
Archaea, and Eukarya \cite{ptc2022, ptc2020, luca2024}.
Specific nucleotide positions are completely invariant in
every sequenced organism. This conservation is enforced by
absolute purifying selection: any mutation in the PTC is
lethal before reproduction.

Therefore: the LUCA's Row~3 product has been produced
continuously in every cell of every lineage since the LUCA
existed. Life has been continuous. The ribosome has
functioned continuously. The Row~3 product has not ceased.
By Definition~2, the LUCA is not extinct. $\square$
\end{proof}

\begin{remark}
The existing literature describes the ribosome as a
``molecular fossil'' and the ``best proxy for LUCA's
biochemistry'' \cite{ptc2022}. This framing is geometrically
imprecise. A fossil is the trace of an extinct organism. The
ribosome is not a trace. It is the active, currently
executing, structurally identical molecular machine that the
LUCA ran. The LUCA is not a historical organism whose traces
persist. It is a currently living identity whose Row~3
expression has been continuous and uninterrupted for
3.8--4.0 billion years. This reframing is not present in
the prior literature.
\end{remark}

\subsection{The ribosome as algebraic-to-geometric transducer}

\begin{definition}[Algebraic-to-Geometric Transducer]
The ribosome is the algebraic-to-geometric transducer of
life: the molecular machine that converts the
one-dimensional algebraic code of genetic sequence
(DNA/mRNA) into three-dimensional geometric structure
(protein fold and function). It is the boundary across
which the informational domain becomes the structural
domain.

Precisely: DNA is to algebra as the ribosome is to geometry.
DNA encodes symbol-to-symbol mappings in a one-dimensional
sequential code. The ribosome converts that algebra into
three-dimensional physical form. Without the ribosome, the
algebra exists but produces no geometry.
\end{definition}

\subsection{The ribosome as topological origin}

\begin{theorem}[The Ribosome as Topological Origin]
The ribosome is the topological origin point of the
structural space of all life.

Every protein in every organism that has ever lived was
produced by the ribosome or by enzymes produced by the
ribosome. The ribosome is not a node in the structural space
of life. It is the operation that generates the space.
Remove any other molecular machine and the space is damaged.
Remove the ribosome and the space does not exist.
\end{theorem}

\subsection{The ribosome as the medium of developmental arrest}

\begin{corollary}[The Ribosome as Medium]
The ribosome is not what is being developmentally arrested
in any protocol in this document. It is the medium on which
every developmental arrest runs.

Every attractor recovered by developmental arrest is
constituted by proteins made by ribosomes. The ribosome
predates every accessible attractor in the evolutionary
table by at least 2.3 billion years (LUCA $\sim$3.8--4.0
Ga; LECA $\sim$1.5 Ga). Therefore when any developmental
arrest recovers any ancestral attractor, the ribosome is
present inside the recovered organism, older than the
organism being recovered, executing the same function it
has executed since before any currently accessible attractor
existed.

The LUCA comes along with every experimental run. It was
never absent.
\end{corollary}

% ══════════════════════════════════════════════════════════════
\section{Falsifiability Summary}
% ══════════════════════════════════════════════════════════���═══

The framework fails at specific, stated points if any of the
following conditions is met. These are Type~B failures in
the sense defined in the repository's
\texttt{falsifiability\_theorem.md}.

\begin{table}[H]
\centering
\caption{Falsifiability conditions. Any one met constitutes
a framework failure at that point and requires explicit
retraction or revision.}
\label{tab:falsify}
\begin{tabularx}{\textwidth}{lX}
\toprule
\textbf{ID} & \textbf{Condition} \\
\midrule
F1 & A coherent alternative trajectory from informational
encoding to a living organism is described that does not
trace the evolutionary path of the lineage. Falsifies the
Reproduction Theorem. \\[0.4em]
F2 & Developmental arrest of a eukaryotic cell at the LECA
stage produces an organism morphologically inconsistent with
LECA-grade characteristics with no structural explanation.
\\[0.4em]
F3 & Arrested LECA-grade organisms from different starting
kingdoms are morphologically distinguishable above chance
by blinded observers. \\[0.4em]
F4 & Developmental arrest at the Ediacaran grade from
different animal lineages produces morphologically distinct
organisms attributable to lineage-specific programs. \\[0.4em]
F5 & A living organism is found whose cells do not contain
ribosomes executing peptide bond formation by the conserved
PTC mechanism. Falsifies the claim that the LUCA is not
extinct. \\[0.4em]
F6 & The developmental hourglass is absent in a major animal
or plant lineage with no structural explanation. Falsifies
the Reproduction Theorem. \\[0.4em]
F7 & The FOXA1/EZH2 ratio fails to order breast cancer
subtypes in a new independent dataset with sufficient sample
size and appropriate platform controls. Falsifies the
Identity Axis Theorem in the BRCA application. \\
\bottomrule
\end{tabularx}
\end{table}

\noindent None of conditions F1--F7 has been met as of
pre-registration date 2026-03-12.

% ══════════════════════════════════════════════════════════════
\section{Priority and Novelty Claims}
% ══════════════════════════════════════════════════════════════

The following claims are made as novel contributions not
present in the existing literature as of 2026-03-12:

\begin{enumerate}[label=\textbf{N\arabic*.}]
  \item \textbf{The Reproduction Theorem as a proved theorem.}
    Proven by elimination of all coherent alternatives. The
    existing literature treats ontogeny-phylogeny
    correspondence as discredited (Haeckel) or approximate
    (evo-devo). Proof by elimination is not in the literature.

  \item \textbf{Developmental arrest as evolutionary access.}
    The arrested organism IS the ancestral attractor --- not
    an analog, not an approximation. Not in the literature.

  \item \textbf{The Inverse Accessibility Theorem.}
    The oldest recoverable organism is the easiest and
    cheapest to recover. Not in the literature.

  \item \textbf{The LECA experiment as the priority
    de-extinction.} The LECA as the deepest accessible
    ancestral attractor, recoverable from any eukaryote,
    cheaper and faster than any de-extinction attempt in
    history. Not proposed in the literature.

  \item \textbf{The cross-kingdom convergence prediction.}
    Human embryo, oak seed, and fungal spore arrested at the
    same developmental depth produce morphologically
    identical organisms. Not in the literature.

  \item \textbf{The LUCA is not extinct.} Geometric
    reframing: the LUCA's Row~3 product has not ceased;
    therefore the LUCA is currently alive. The ``molecular
    fossil'' framing is geometrically incorrect. Not in the
    literature.

  \item \textbf{The ribosome as topological origin.}
    The ribosome as the origin point of the structural space
    of life --- the algebraic-to-geometric transducer. Not
    in the literature.

  \item \textbf{The ribosome as medium of all developmental
    arrest.} The ribosome is present inside every recovered
    ancestral organism, older than every attractor it will
    ever be found in. The LUCA comes along with every
    experimental run. Not in the literature.
\end{enumerate}

% ══════════════════════════════════════════════════════════════
\section{Locked Statement}
% ══════════════════════════════════════════════════════════════

\begin{center}
\colorbox{warnbox}{
\begin{minipage}{0.88\textwidth}
\vspace{0.6em}
\textbf{This record is locked as of 2026-03-12.}

\medskip
All predictions in this document were derived from first
principles using the triadic invariant and the Reproduction
Theorem. All predictions were locked before any experimental
confirmation. All prior predictions in this framework
(including the FOXA1/EZH2 ratio across $\sim$7,500 patients
in seven independent datasets) were confirmed in the correct
direction with zero biological contradictions.

\medskip
\textbf{Reproduction Theorem:} Proved by elimination of all
coherent alternative trajectories. Open for 150 years. No
alternative has been described.

\medskip
\textbf{LECA experiment:} \$200--400 per run. Any eukaryote.
Any standard cell biology laboratory. Days to result. The
oldest recoverable organism. The easiest to recover.

\medskip
\textbf{Cross-Kingdom Convergence Prediction:} Human embryo
and oak seed and fungal spore, arrested at the same
developmental depth, will produce morphologically identical
organisms. Geometrically required. Not yet experimentally
confirmed. Staked here.

\medskip
\textbf{LUCA:} Not extinct. The ribosome is the LUCA. The
LUCA is in every cell of every organism alive on Earth
right now, executing the same function it has executed for
3.8--4.0 billion years without interruption. Not a fossil.
The organism itself. Still running. Inside everything.

\medskip
\textit{``Nothing was lost. All of it is running right now.
All of it is accessible. The geometry says so.''}

\vspace{0.4em}
\hfill --- Eric Robert Lawson, March 12, 2026
\vspace{0.4em}
\end{minipage}
}
\end{center}

% ══════════════════════════════════════════════════════════════
\section*{Repository and Reproducibility}
\addcontentsline{toc}{section}{Repository and Reproducibility}
% ══════════════════════════════════════════════════════════════

All derivation records, protocols, reasoning artifacts, and
pre-registration documents are publicly available at:

\begin{center}
\href{https://github.com/Eric-Robert-Lawson/attractor-oncology}
{\texttt{https://github.com/Eric-Robert-Lawson/attractor-oncology}}\\[0.3em]
\small Repository ID: 1172048282 \textbar{} MIT License
\textbar{} Python \textbar{} Public
\end{center}

\noindent Key files relevant to this pre-registration:

\begin{center}
\begin{tabular}{ll}
\toprule
\textbf{File} & \textbf{Location in repository} \\
\midrule
FOXA1/EZH2 master artifact
  & \texttt{Cancer\_Research/BRCA/.../FOXA1\_EZH2\_BCRA\_Ratio.tex} \\
Reproduction Theorem
  & \texttt{THE\_REPRODUCTION\_THEOREM.md} \\
Evolutionary Table
  & \texttt{THE\_EVOLUTIONARY\_TABLE.md} \\
Ediacaran Convergence Prediction
  & \texttt{THE\_EDIACARAN\_CONVERGENCE\_PREDICTION.md} \\
Dickinsonia Protocol
  & \texttt{THE\_DICKINSONIA\_PROTOCOL.md} \\
LECA Resurrection Protocol
  & \texttt{THE\_LECA\_RESURRECTION\_PROTOCOL.md} \\
Plant Ediacaran and LECA Attractor
  & \texttt{THE\_PLANT\_EDIACARAN\_AND\_THE\_LECA\_ATTRACTOR.md} \\
Complete Ancestral Map
  & \texttt{THE\_COMPLETE\_ANCESTRAL\_MAP.md} \\
Inverse Accessibility Theorem
  & \texttt{THE\_INVERSE\_ACCESSIBILITY\_THEOREM.md} \\
LUCA Is Not Extinct
  & \texttt{THE\_LUCA\_IS\_NOT\_EXTINCT.md} \\
Ribosome as Topological Origin
  & \texttt{THE\_RIBOSOME\_AS\_TOPOLOGICAL\_ORIGIN.md} \\
Ribosome as Medium
  & \texttt{THE\_RIBOSOME\_AS\_MEDIUM.md} \\
Falsifiability Theorem
  & \texttt{falsifiability\_theorem.md} \\
On the Author and the Geometry
  & \texttt{ON\_THE\_AUTHOR\_AND\_THE\_GEOMETRY.md} \\
\bottomrule
\end{tabular}
\end{center}

% ══════════════════════════════════════════════════════════════
\section*{Competing Interests}
\addcontentsline{toc}{section}{Competing Interests}
% ══════════════════════════════════════════════════════════════

The author has no competing interests to declare. No
institutional funding was received. No conflicts of interest
exist. All data used is publicly available. All code is
open source under the MIT licence.

% ══════════════════════════════════════════════════════════════
\section*{Author Statement}
\addcontentsline{toc}{section}{Author Statement}
% ══════════════════════════════════════════════════════════════

This work was conducted independently by Eric Robert Lawson
under the OrganismCore research programme. The author has no
formal training in biology and no institutional affiliation.
The framework was derived from first principles using
geometric reasoning. No laboratory was used. No collaboration
was involved. The geometry was the method. The pre-specified
prediction record and the zero Type~B failure rate across all
domains constitute the epistemic foundation of the claims.

% ══════════════════════════════════════════════════════════════
\section*{Contact}
\addcontentsline{toc}{section}{Contact}
% ══════════════════════════════════════════════════════════════

\begin{tabular}{ll}
Author        & Eric Robert Lawson \\
Organisation  & OrganismCore \\
Email         &
  \href{mailto:OrganismCore@proton.me}
       {OrganismCore@proton.me} \\
ORCID         &
  \href{https://orcid.org/0009-0002-0414-6544}
       {https://orcid.org/0009-0002-0414-6544} \\
Repository    &
  \href{https://github.com/Eric-Robert-Lawson/attractor-oncology}
       {github.com/Eric-Robert-Lawson/attractor-oncology} \\
Date          & 2026-03-12 \\
\end{tabular}

% ── BIBLIOGRAPHY ──────────────────────────────────────────────
\newpage
\bibliographystyle{plainnat}
\begin{thebibliography}{99}

\bibitem{kalinka2010}
Kalinka, A.T., Varga, K.M., Gerrard, D.T., Preibisch, S.,
Corcoran, D.L., Jarrells, J., Ohler, U., Bergman, C.M.,
\& Tomancak, P. (2010).
Gene expression divergence recapitulates the developmental
hourglass model.
\textit{Nature}, 468(7325), 811--814.
\doi{10.1038/nature09634}

\bibitem{irie2011}
Irie, N., \& Kuratani, S. (2011).
Comparative transcriptome analysis reveals vertebrate
phylotypic period during organogenesis.
\textit{Nature Communications}, 2, 248.
\doi{10.1038/ncomms1248}

\bibitem{plant_hourglass2024}
[Authors TBC]. (2024).
Multiplexed transcriptomic analyses of the plant embryonic
hourglass.
\textit{Nature Communications}.
\url{https://www.nature.com/articles/s41467-024-55803-9}

\bibitem{line1_2024}
[Authors TBC]. (2024).
LINE-1 elements prevent developmental regression and lock
forward development downstream of global genome silencing.
\textit{Nature}.
[LINE-1 forward-lock mechanism, 2024.]

\bibitem{ptc2022}
Rozov, A., \& Yusupov, M. (2022).
The peptidyl transferase center: a window to the past.
\textit{Microbiology and Molecular Biology Reviews},
86(1).
\doi{10.1128/mmbr.00104-21}

\bibitem{ptc2020}
Petrov, A.S., et al. (2020).
The ancient history of peptidyl transferase center
formation as told by conservation and diversification
of its components.
\textit{Life}, 10(8), 134.
\doi{10.3390/life10080134}

\bibitem{luca2024}
Moody, E.R.R., et al. (2024).
Metabolism, genome and age of the last universal common
ancestor.
\textit{Nature}, 628, 545--550.
\doi{10.1038/s41559-024-02474-w}

\bibitem{srivastava2008}
Srivastava, M., et al. (2008).
The \textit{Trichoplax} genome and the nature of placozoans.
\textit{Nature}, 454(7207), 955--960.
\doi{10.1038/nature07191}

\bibitem{bobrovskiy2018}
Bobrovskiy, I., et al. (2018).
Ancient steroids establish the Ediacaran fossil
\textit{Dickinsonia} as one of the earliest animals.
\textit{Science}, 361(6408), 1246--1249.
\doi{10.1126/science.aat7228}

\bibitem{schade2024}
Schade, A.E., et al. (2024).
PRC2 is the convergence node in TNBC silencing
FOXA1/GATA3/ESR1.
\textit{Nature}.
[Independent confirmation of EZH2 as convergence hub
in TNBC; confirms FOXA1/EZH2 mechanistic basis.]

\bibitem{toska2017}
Toska, E., et al. (2017).
PI3K pathway regulates ER-dependent transcription in
breast cancer through the epigenetic regulator KMT2D.
\textit{Nature Medicine}, 23(10), 1170--1178.
\doi{10.1038/nm.4402}

\bibitem{wu2021}
Wu, S.Z., et al. (2021).
A single-cell and spatially resolved atlas of human
breast cancers.
\textit{Nature Genetics}, 53(9), 1334--1347.
\doi{10.1038/s41588-021-00911-1}

\bibitem{ramesh2005}
Ramesh, M.A., Malik, S.B., \& Logsdon, J.M. (2005).
A phylogenomic inventory of meiotic genes: evidence for
sex in \textit{Giardia} and an early eukaryotic origin
of meiosis.
\textit{Current Biology}, 15(2), 185--191.
\doi{10.1016/j.cub.2005.01.003}

\end{thebibliography}

% ═���════════════════════════════════════════════════════════════
\end{document}
